\documentclass[12pt]{article}

\usepackage{eurosym}
\usepackage{amsfonts,amsmath,amssymb,graphicx,setspace,natbib,bbm,subcaption}
\usepackage[usenames, dvipsnames]{xcolor}
\usepackage[colorlinks=true,linkcolor=blue,citecolor=Mahogany,urlcolor=blue]{hyperref}
\usepackage[top=1in, bottom=1in, left=1in, right=1in]{geometry}
\usepackage[normalem]{ulem}
\usepackage{xcolor}
\usepackage{parskip}

\newtheorem{theorem}{Theorem}
\newtheorem{proposition}{Proposition}
\newtheorem{corollary}{Corollary}
\newtheorem{lemma}{Lemma}
\newtheorem{definition}{Definition}
\newtheorem{assumption}{Assumption}
\newtheorem{conjecture}{Conjecture}
\newtheorem{observation}{Observation}
\newtheorem{condition}{Condition}
\newtheorem{property}{Property}
\newtheorem{claim}{Claim}
\newtheorem{remark}{Remark}
\newenvironment{proof}[1][Proof]{\noindent\textbf{#1.} }{\ \rule{0.5em}{0.5em}}
\newenvironment{question}[1][Question]{\noindent\textbf{#1.} }{\ \rule{0.5em}{0.5em}}

\newcommand\E{{\rm E}}
\newcommand{\blue}[1]{\textcolor{blue}{#1}}

\definecolor{darkgreen}{rgb}{0,0.5,0}
\newcommand{\nic}[1]{\noindent\textcolor{darkgreen}{[Nic: #1]}}

\allowdisplaybreaks
\onehalfspacing

\begin{document}

\section*{Organizing hypotheses}

    Goal: explore and rule out competing hypotheses for why young people in the U.S. are more likely to support redistribution (either this generation specifically or all generations when they are young). 

    Framework: 
    \begin{equation*}
        U = f(\text{consumption + leisure + values/morals})
    \end{equation*}
    There are \textbf{preferences} which shift the weights assigned to consumption, leisure, and morals in the utility function. There are \textbf{beliefs} which shift the values of consumption, leisure, and morals. There are \textbf{behavioral biases} which shift one's perceived preferences--they think they know what their utility function looks like, but they don't. 

    \subsection*{Preferences}

    Consumption preferences:
    \begin{itemize}
        \item Uncertain future income, so they want to reduce risk.
        \begin{itemize}
            \item Min max
        \end{itemize}
        \item High discount rate, so they want to maximize current income, which is low.
        \item Young people care more/less about economic growth or materialistic things
        \begin{itemize}
            \item Or they're materialistic, but it's limited. They don't see much extra value in a very high income. Related to the concept of how much value a society places on financial success/social standing.
        \end{itemize}
        \item Older people care more about having high incomes because of houses, children, etc.
        \begin{itemize}
            \item Young people might think right now that more income is a luxury and not a necessity.
            \item Or, they might not understand how much money you need to live on when you're older.
        \end{itemize}
    \end{itemize}

    Leisure preferences:
    \begin{itemize}
        \item Young people prefer work-life balance, or find hard work painful, or place more value on activities outside work.
    \end{itemize}

    Values/moral preferences:
        \begin{itemize}
        \item Individualism. Individualists are those who place individual needs before society. Similar to general selfishness.
        \begin{itemize}
            \item Relates to materialism.
            \item Relates to freedom
            \item Relates to compassion/care for others.
        \end{itemize}
        \item Ingroup vs. outgroup (minority vs. majority) preferences
        \begin{itemize}
            \item Role of economic shocks
            \item Usually in reference to race
            % \item Prejudice against lower class people which changes as you age
        \end{itemize}
        \item Openness to new experience is correlated with liberal voting. 
        \begin{itemize}
            \item Younger people may be more open to experience.
            \item Or, young people today may be more open to experience because of education or exposure to the world through media.
        \end{itemize} 
        \item Universalist vs. communalist/particularist. Universalists care about concepts such as rights and fairness that apply to all people, while communalists favor their own group and emphasize values such as loyalty and tradition. Communalists support nationalist candidates.
        \begin{itemize}
            \item Relates to fairness
        \end{itemize}
            
    \end{itemize}

    \subsection*{Beliefs}
    
    Consumption beliefs:
    \begin{itemize}
        \item Incentive effects
        \begin{itemize}
            \item Learning from family's mobility experience
            \item Learning from personal mobility experience
            \item Learning from recent history and recent generations' experience
            \item Media
            \item Optimism
            \begin{itemize}
                \item Relates to openness to experience
                \item Young people haven't had bad experience with government, so they overstate its effectiveness and support it more
            \end{itemize}
            % \item Zero-sum game
        \end{itemize}
    \end{itemize}

    \subsection*{Biases}

    Behavioral biases:
    \begin{itemize}
        \item Young people don't anticipate how much taxes will hurt
        \begin{itemize}
            \item Relates to ``lack of real-world experience'', which also informs some other beliefs.
        \end{itemize}
        \item Myopic
        \begin{itemize}
            \item Relates to having a high discount factor, but in this case, you're not consciously discounting your future income; you're falsely estimating it or not even considering it.
        \end{itemize}
    \end{itemize}

    


\end{document}